\section{Discussion, limitations and next steps}\label{sec:discussion}

\paragraph{What carries over to a production fulfilment network.}
The pipeline is modular in the same way production planning stacks are:
\begin{itemize}
  \item forecasters produce quantiles;
  \item a scenario layer adds dependence;
  \item a stochastic optimiser commits capacity;
  \item a fast LP-dual controller assigns orders against that capacity.
\end{itemize}
Each layer can be replaced without touching the others, and each comes with its own evaluation: proper scores, joint-calibration scores, statistical optimality gaps and regret to a clairvoyant benchmark.
The \emph{decision-focused} evaluation of \cref{sec:value-forecast} is the piece we would most want in production.
It prices forecast improvements in euros per day, which is the currency in which forecasting and planning teams should negotiate their roadmaps.

\paragraph{Limitations.}
\begin{enumerate}
  \item \textbf{Synthetic data.} All demand is simulated. The generative model was designed to contain the phenomena that make the real problem hard: events, holidays, promotions, weather, correlated shocks and overdispersion. It nevertheless remains a model, and real data would contain regime changes, censored demand (stock-outs) and assortment effects that it lacks. The upside of this choice is that it enables the oracle comparisons that are impossible with real data.
  \item \textbf{Stylised network and costs.} Sites, capacities and costs are plausible, not real. The structural conclusions (pooling, the role of recourse flexibility, the event-day risk of naive forecasts) are robust to the exact numbers; the euro amounts are not.
  \item \textbf{Single-unit shipments.} Multi-item orders introduce split-shipment costs and the correlated rounding problem of \citet{jasin2015,ma2023}. The online simulator would need order-level decisions.
  \item \textbf{Single-day horizon.} Unshipped units are penalised but not carried over. A rolling multi-day model, with backlog dynamics and next-day promise re-planning, would turn the two-stage model into a multi-stage one.
  \item \textbf{Exogenous inventory.} Inventory placement, which ultimately determines which FCs \emph{can} ship on time, is taken as given. Joint inventory placement and capacity planning is a natural (and harder) extension.
  \item \textbf{Network size.} The decomposition is fast for the Iberian network but not yet for a network twice its size (\cref{sec:computation}): the serial master MIP is the bottleneck.
\end{enumerate}

\paragraph{Next steps.}
\begin{itemize}
  \item \emph{Forecasting:} fine-tune the foundation model with covariates; hierarchical reconciliation of the zone-level forecasts to the national level; a copula with a factor structure (national, regional and category factors) instead of a shrunk full correlation.
  \item \emph{Optimisation:} branch-and-Benders-cut in a single search tree with lazy constraints; Magnanti--Wong or Papadakos Pareto-optimal cuts; a risk-averse objective (CVaR of the daily cost) through the standard Rockafellar--Uryasev reformulation, which keeps the Benders structure.
  \item \emph{Learning:} a learned surrogate of the recourse function, trained on the millions of subproblem solutions the decomposition produces, to warm-start the master or screen scenarios.
  \item \emph{Online control:} bid prices from a stochastic rather than a deterministic fluid LP, and multi-armed-bandit tuning of the re-solve schedule against live regret.
\end{itemize}
